\documentclass[11pt]{article}
\usepackage[T1]{fontenc}
\usepackage[utf8]{inputenc}
\usepackage{amsmath,amssymb,amsthm,booktabs,hyperref}
\usepackage[margin=1in]{geometry}
\hypersetup{hidelinks}
\newtheorem{theorem}{Theorem}[section]
\newtheorem{lemma}[theorem]{Lemma}
\newtheorem{proposition}[theorem]{Proposition}
\newtheorem{corollary}[theorem]{Corollary}
\theoremstyle{definition}
\newtheorem{definition}[theorem]{Definition}
\newtheorem{remark}[theorem]{Remark}
\title{Source-indexed information repair:\\stable representations, lawful acquisition, and online realization}
\author{Observational SUCC research program\\Technical manuscript; academic priority unverified}
\date{19 October 2026}
\begin{document}
\maketitle
\begin{abstract}
We distinguish three finite mathematical problems: autonomously retaining
future-task information, lawfully acquiring an original source label through
possibly destructive probes, and realizing abstract requests online. Stable
partition refinement gives a canonical deterministic repair and a sharp
fixed-width auxiliary-memory bound. For acquisition, a source-indexed belief
construction yields exact rational minimax values with locally checkable
cost and termination certificates. The minimal retained advice at a given
probe budget is characterized by weak coloring of the hypergraph of minimal
infeasible source sets. A three-source example has all pairs identifiable at
unit cost but no joint identification strategy. The general ingredients are
classical; the contribution is a precise integration with explicit resource,
provenance, and information boundaries, not a claim of established novelty.
\end{abstract}

\section{Scope and mathematical interfaces}
All carriers in the constructive results are complete finite declared models.
An original source label records provenance as a mathematical parameter; it
does not authenticate physical reality. All costs are nonnegative rationals.
There are no probabilistic, asymptotic, physical-energy, or Riemann-hypothesis
conclusions.

Our baseline is PR179, commit
\texttt{81684b4b8956f9f5da56c376ba7e4339277eb665}.
Sibling PR178, commit
\texttt{e42096ec0186228dd0f748767599e470f2ca2e1d}, was inspected separately.
The sibling's invariant private source tag differs from the baseline's
preserved source output; their failed-run journals also differ. We reuse
their common local kernel-congruence requirement, not an identification of
their full semantics. In particular, acquisition below has no refusal oracle:
all probes must be lawful throughout the current uncertainty set.

\section{The least autonomous repair}
Fix a nonempty finite state space $X$, output $o:X\to O$, and partial
deterministic actions $\delta_a:D_a\to X$ with costs
$c_a:D_a\to\mathbb Q_{\ge0}$. Any immutable source identity or designated
proof-task result that must be retained is included in $o$. Let $q:X\to Y$
be an existing representation. An equivalence is \emph{stable} when it
preserves output, enabledness, cost, and successor equivalence for every
action. Let $R_*$ be the greatest stable equivalence contained in $\ker q$.

\begin{theorem}[Classical stable refinement and factorization]
$R_*$ exists and is obtained by finite partition refinement starting from
$\ker(q,o)$. There are at most $|X|-|(q,o)(X)|$ strict refinements. If
$r:X\twoheadrightarrow Z$ is any autonomous sufficient repair retaining $q$,
there is a unique map $f:Z\to X/R_*$ such that $\pi_*=f\circ r$.
\end{theorem}
\begin{proof}
Refine by each action's unavailability or pair of cost and previous successor
class. Agreement through $k$ actions is exactly the $k$th partition, by
induction. A plateau is stable. Every strict refinement increases the number
of blocks, giving termination. Any stable equivalence refining $(q,o)$
refines every stage by induction, hence refines the terminal partition.
The kernel of $r$ is such an equivalence. Therefore $f(r(x))=[x]_{R_*}$ is
well-defined and unique by surjectivity.
\end{proof}

This is enriched Moore/Nerode minimization, within classical coalgebraic
semantics \cite{Moore,Rutten}. It is not merely the intersection of $\ker q$
with the old behavioral equivalence: future values of $q$ must be preserved
as well. For example, let all old outputs be constant and let the only action
send $u\mapsto u$, $v\mapsto w$, $w\mapsto w$, at zero cost. If
$q(u)=q(v)\ne q(w)$, the old behavior is constant, but the repair separates
all three states.

\begin{theorem}[Sharp final auxiliary alphabet]
Put $K=\max_y|q^{-1}(y)/R_*|$. The smallest alphabet for an auxiliary label
$\alpha:X\to\{0,\ldots,k-1\}$ making $(q,\alpha)$ an autonomous sufficient
repair has size $K$. Its smallest fixed binary width is $\lceil\log_2K\rceil$.
\end{theorem}
\begin{proof}
Any repair kernel lies inside $R_*$. Within one $q$-fiber, distinct $R_*$
classes therefore need distinct auxiliary labels. Conversely, enumerate the
classes separately in each fiber and reuse the same alphabet between fibers.
Then $\ker(q,\alpha)=R_*$. Binary capacity gives the ceiling formula.
\end{proof}

The encoder may inspect the source state, while the decoder knows $q$.
This is an existence and retention bound, not a construction of an observation
instrument. The supplemental label need not update autonomously without $q$.
No smaller code is obtained by supposing an unperformed acquisition occurred.

\section{Acquiring an original label through destructive probes}
Let $X_0$ be the original source set, $Z$ a finite carrier, $i:X_0\to Z$ an
initialization, and $t:X_0\to T$ the target original repair class. A probe
maps its domain in $Z$ to receipt, exact cost, and new carrier. A belief is
a nonempty set $B\subseteq X_0\times Z$ with at most one carrier for each
original source. A probe is lawful only if defined at every carrier in $B$.
For a receipt/cost outcome $h=(\rho,c)$, its successor is
\[
 B_{a,h}=\{(x,d_a(z)):(x,z)\in B,\ (\rho_a(z),c_a(z))=h\}.
\]
The original coordinate remains fixed even when current carriers merge.
Success means all originals in the terminal belief have the same $t$-value.
Policies must terminate on every possible original. Their cost is the
maximum total charge over those originals.

\begin{theorem}[Exact terminating acquisition cost]
Let the reachable belief arena, stopped at success, have $N$ states. Initialize
$V_0=0$ at success and $V_0=\infty$ elsewhere. At nonsuccessful beliefs set
\[
 V_{n+1}(B)=\min\left(V_n(B),\min_{a\text{ lawful}}
       \max_{h:B_{a,h}\ne\varnothing}\{c(h)+V_n(B_{a,h})\}\right).
\]
Then $V_N$ is the exact optimum among terminating policies. A finite optimum
is attained and rational.
\end{theorem}
\begin{proof}
Induction identifies $V_n$ with optimal identification trees of height at most
$n$. If a finite successful tree repeats a belief on a branch, replace its
earlier occurrence by the later subtree. The source/current belief is equal,
so the substituted policy is lawful and correct. Deleted costs are
nonnegative; no original source pays more. Repeated deletion leaves height
at most $N$. There are finitely many trees of this bounded height, proving
attainment and the stated recurrence bound.
\end{proof}

\begin{proposition}[Local optimality certificate]
An exact value certificate assigns $L(B)\in\mathbb Q_{\ge0}\cup\{\infty\}$.
Success states have value zero and stop. Every lawful action obeys
\[
 L(B)\le\max_h\{c(h)+L(B_{a,h})\}.
\]
At each finite nonsuccessful state a selected action obeys the reverse
inequality, has finite-valued children, and strictly decreases a supplied
nonnegative integer rank along every outcome. These conditions imply $L=V$.
\end{proposition}
\begin{proof}
Induction on any finite policy tree proves that $L$ lower-bounds its worst
cost. This also excludes a finite successful policy at an infinity-labelled
state. Induction on the selected policy's decreasing rank proves termination
and an upper cost bound $L$. The two inequalities give equality.

Such ranks exist for the optimal values. Call an action optimal when its
cost-plus-optimal-child maximum is at most the parent's finite value.
Form the goal attractor using only these actions. If finite-valued states
remain outside, choose their minimum value $m$, then a state of value $m$
with a shortest optimal tree. Its root action is optimal, so some child is
still outside the attractor. Nonnegative cost and minimality force that
child's value to be $m$ and its entering cost to be zero. Its subtree has
cost at most $m$, hence is optimal and shorter, a contradiction. Attractor
ranks therefore supply one proper positional optimal policy. This argument
does not assume that every branch of a root-optimal tree is locally optimal.
\end{proof}

A zero-cost nonidentifying loop satisfies the spurious Bellman equation
$v=0+v$ with $v=0$; it fails the rank condition. Thus an arbitrary fixed point
is not a termination certificate. The synthesis makes no claim that
destructive acquisition restores the carrier for later task execution.
That conclusion needs a separate source-preserving restoration or simulation
map. Full transcript/controller memory is available during acquisition.

\section{The hypergraph of acquisition obstructions}
For budget $C$ and original observation $y$, let $\mathcal F_C(y)$ be the
nonempty subsets of $q^{-1}(y)$ identifiable within cost $C$. Restricting a
policy to fewer originals preserves lawfulness and cannot increase worst-case
cost. Consequently $\mathcal F_C(y)$ is hereditary and contains all
singletons. Let $\mathcal H_C(y)$ be its minimal infeasible subsets.

\begin{lemma}[Finite hereditary families]\label{lem:hereditary}
A finite set is feasible in a hereditary family exactly when it contains no
minimal infeasible subset. The minimum number of feasible blocks partitioning
the vertex set equals the weak chromatic number of its minimal-infeasible
hypergraph, where no hyperedge may be monochromatic.
\end{lemma}
\begin{proof}
Every infeasible finite set contains a minimal one; a feasible set cannot
contain an infeasible subset. A coloring is therefore valid exactly when
each color class is feasible.
\end{proof}

\begin{theorem}[Exact prior-advice/acquisition tradeoff]
Suppose $q(x)$ and an auxiliary source label $\alpha(x)$ are already retained
before probing. The minimum common advice alphabet permitting cost-$C$
identification is
\[
 K_{\rm prior}(C)=\max_y\chi(\mathcal H_C(y)),\qquad
 b_{\rm prior}(C)=\lceil\log_2K_{\rm prior}(C)\rceil.
\]
\end{theorem}
\begin{proof}
After observing $(y,\alpha)$, the whole joint fiber must be identifiable.
Thus every advice color class must be feasible. Conversely, use a feasible
policy for each color class of a valid coloring. Apply
Lemma~\ref{lem:hereditary} within each $q$-fiber and reuse colors between
fibers. The binary formula follows from alphabet capacity.
\end{proof}

Advice may split inside a final repair class. It must have been acquired or
retained earlier; this theorem does not provide a free source oracle. Prior
advice width, final autonomous repair width, and temporary controller memory
are three distinct quantities. Budget breakpoints are among the finitely
many exact initial-subset values. Exhaustive finite synthesis gives a precise
frontier, without a polynomial-time claim.

\subsection*{A three-source obstruction and sharp frontier}
Originals $x,y,z$ share $q$ and have distinct target labels. A unit-cost probe
$a$ distinguishes $x$ but sends $y,z$ to the same carrier $u$ with the same
receipt. A unit-cost probe $b$ distinguishes $z$ but merges $x,y$ into carrier
$v$. At $u,v$ every probe loops with the same receipt. Each pair is identifiable
at cost 1, but every first action from the triple has an irreversibly
indistinguishable branch. No joint identification policy exists.

At budget 1 the unique minimal forbidden set is $\{x,y,z\}$. Its pairwise
conflict graph has no edges, but its weak chromatic number is 2. Add a
nondestructive identifying probe of cost 2 at $x,y,z$. The exact frontier is
\[
 (b_{\rm prior},C)=(2,0),\quad(1,1),\quad(0,2).
\]
This proves that pairwise minimum distinguishing costs do not determine
joint acquisition feasibility. Forgetting the original coordinate would
falsely declare the merged carrier resolved and change the task itself.

% Section fragment for paper_a_source_information.tex.
% Uses the parent's amsmath/amssymb and theorem/proposition environments.
\section{Online source realization and public command interfaces}
\label{online-realization}

Offline history existence and online control have different quantifiers.
For a finite source with private states $X$, reports $q:X\to Y$, and lawful
labelled transitions $x\xrightarrow a x'$, fix an abstract transition relation
on $Y$. Original provenance is a map $\sigma:X\to\Sigma$ preserved by every
source edge. Current carrier and epoch may change inside $x$. Neither this
metadata nor a mathematical observation refinement is automatically public.

In the perfect-information realization game, the environment requests an
abstract transition $q(x)\xrightarrow a y'$. The resolver then chooses a
matching private transition $x\xrightarrow a x'$ with $q(x')=y'$, before seeing
the next request. An unanswered request is a loss. Absence of any abstract
request means safe termination, not satisfaction of a progress objective.

\subsection{Exact winning regions and finite certificates}
\label{online-certificates}

It is useful to allow a further adversarial outcome after the resolver's
command. Let $V$ be finite; let $\operatorname{Req}(v)$ be its requests,
$\operatorname{Resp}(v,r)$ its commands, and
$\varnothing\ne\operatorname{Out}(v,r,c)\subseteq V$ their complete outcome
sets. Empty response sets express refusal. Define
\[
 F(Z)=\{v\in V:\ \forall r\in\operatorname{Req}(v)\ \exists c\in
 \operatorname{Resp}(v,r),\ \operatorname{Out}(v,r,c)\subseteq Z\}.
\]
The source-realization game is the special case in which each selected private
successor is a response with a singleton outcome set.

\begin{theorem}[Finite sure-safety certificate]
\label{online-safety}
The greatest fixed point $W=\nu F$ is exactly the sure-winning region, and
positional resolver strategies suffice. Its complement has ranks
$\rho:V\setminus W\to\{1,\ldots,|V|\}$, a chosen request at every losing
position, and for every response a chosen lower-rank outcome. These choices
force refusal within at most $\rho(v)$ requests from $v$.

A positive response table closed in $W$, together with a ranked negative table
on its complement, certifies the exact winning region independently of the
algorithm that found it.
\end{theorem}

\emph{Proof.}
The map $F$ is monotone. Iterating $Z_0=V$, $Z_{k+1}=F(Z_k)$ gives a decreasing
sequence stabilizing after at most $|V|$ strict deletions. Every postfixed set
is contained in every $Z_k$, so its limit is $\nu F$. Choose in this limit one
response whose complete outcome set remains in the limit for every request.
Induction on play length proves sure safety.

Give a state first removed at stage $k$ rank $k$. Its removal supplies a
request for which every response has an outcome outside $Z_{k-1}$, hence of
smaller rank. At rank one there is no response. Following these chosen outcomes
strictly decreases a positive integer after every completed round and defeats
any history-dependent resolver. The same arguments verify any supplied pair
of positive and ranked-negative tables directly. Their regions cannot overlap;
if they partition $V$, they identify the winning region exactly.\hfill$\square$

This is classical safety-game and alternating-simulation mathematics; compare
Alur, Henzinger, Kupferman and Vardi, \emph{Alternating Refinement Relations}
(CONCUR 1998). Its role here is to make control ownership and refusal explicit.

For example, take states $p,u,v$ with one report and initial state $p$.
Either action $a,b$ sends $p$ to either $u,v$. Only $a$ is available at $u$,
and only $b$ at $v$; either enabled action again permits both successors.
Every finite word has a private lift, since its support becomes $\{u,v\}$
after the first letter and stays there. Yet online ranks are
$\rho(p)=2$, $\rho(u)=\rho(v)=1$, so there is no online resolver. This is a
finite history-determinism obstruction, not a new separation from automata
theory; compare Boker and Lehtinen, \emph{Token Games and History-Deterministic
Quantitative Automata}, arXiv:2110.14308.

\subsection{Public observation requires a command interface}
\label{online-public-command-interface}

When $x$ is hidden, ``choose a successor'' is not a public action specification.
Declare instead commands $c\in C$ and complete implementation rows
$T_{a,c}(x)\subseteq X$, each contained in the lawful $a$-successors of $x$.
Nature chooses their outcomes adversarially. The public observer sees
$h:X\to O$, where $h$ retains $q$, and has perfect recall of signals, requests,
and its own commands. It does not know the original label unless that label
is available through $h$. Beliefs may contain multiple original labels.

At a nonempty homogeneous belief $B$, the command $c$ answers request $(a,y')$
only if
\[
 \forall x\in B:\quad
 \varnothing\ne T_{a,c}(x)\subseteq q^{-1}(y').
\]
All actual outcomes must be checked before conditioning: discarding outputs
with the wrong report would introduce postselection. After a legal command
and observed signal $o$, update by
\[
 B_o=\left(\bigcup_{x\in B}T_{a,c}(x)\right)\cap h^{-1}(o),
 \qquad B_o\ne\varnothing.
\]
The belief game has these sets as positions and all nonempty $B_o$ as
adversarial command outcomes. It has at most
$\sum_o(2^{|h^{-1}(o)|}-1)$ positions.

\begin{theorem}[Public sure-safety reduction]
\label{online-public-safety}
A pure public-history strategy is surely safe from $B$ if and only if $B$ is
winning in the belief game. A strategy depending on the current belief and
request suffices, although the current report alone may be insufficient.
\end{theorem}

\emph{Proof.}
Induction over a public history shows that the belief is exactly the endpoints
of all private histories compatible with it. A command must work at each such
endpoint, and its possible next signals yield exactly the displayed update.
Thus a winning belief strategy yields a lawful public strategy. Conversely,
fix a public strategy at a losing belief. The ranked certificate supplies a
finite sequence of requests and successor signals leading to a command which
is unavailable or has a wrong-report outcome at some compatible endpoint.
The whole finite signal history has one compatible private history by the
same support induction. Choosing that witness gives a concrete failure of
sure safety. This reasoning never switches the realized original source
between successive steps.\hfill$\square$

This is the safety specialization of the knowledge construction of Chatterjee,
Doyen, Henzinger and Raskin, \emph{Algorithms for Omega-Regular Games with
Imperfect Information}, LMCS 3(3:4), 2007, Section 3.1. It makes no claim
about almost-sure objectives or stochastic belief sufficiency.

\begin{proposition}[The source relation does not identify public control]
\label{online-command-nonidentifiability}
The same private transition relation, initial uncertainty and public reports
can have opposite public online winners under different command interfaces.
\end{proposition}

\emph{Proof.}
Take indistinguishable initial states $0,1$, a terminal state $d$, and edges
$0\xrightarrow a d$, $1\xrightarrow a d$. One command implementing both rows
is surely successful. Two commands, one available only at $0$ and one only
at $1$, have the same union of private edges but neither is surely legal at
belief $\{0,1\}$.\hfill$\square$

\subsection{Probe repair: a structural no-go and a hypergraph criterion}
\label{online-probe-repair}

\begin{proposition}[Information cannot repair structural losing states]
\label{online-structural-no-go}
If $x$ is losing in the perfect-information private-successor game, no truthful
nondestructive probes and public memory can make it surely winning while
keeping the source edges and requested abstract relation fixed.
\end{proposition}

\emph{Proof.}
A perfect-information private resolver can emulate any such public controller
and choose a lawful outcome of each selected command. Sure success of the
public controller would contradict the private losing rank certificate.
Equivalently, observing more information does not create an edge that escapes
the certificate's rank decrease.\hfill$\square$

The converse fails by Proposition~\ref{online-command-nonidentifiability}.
An information defect and a structural control defect require different
repairs. A state-changing probe is outside this no-go's assumptions.

For an exact positive result, consider a one-step terminal control problem
with one current report and one request. Let $E_x\subseteq C$ be the commands
that successfully complete the task at private state $x$. Let $\mathcal H$
contain precisely the inclusion-minimal nonempty sets $H\subseteq X$ with
$\bigcap_{x\in H}E_x=\varnothing$. Supply a finite menu of total truthful
nondestructive probes $p:X\to O_p$, each with an exact nonnegative cost $c_p$.
A fixed bundle is acquired before the request, and its outputs are retained.

\begin{theorem}[Minimal fixed-bundle command repair]
\label{online-probe-hitting}
A probe bundle $S$ permits a surely successful command if and only if every
$H\in\mathcal H$ is split by some $p\in S$. Hence minimum acquisition cost
is exactly the weighted hitting-set optimum whose probe $p$ covers the
hyperedges on which it is nonconstant. A hyperedge split by no available
probe certifies impossibility. A singleton hyperedge cannot be repaired by
any nondestructive observation.
\end{theorem}

\emph{Proof.}
Apply Lemma~\ref{lem:hereditary} to the hereditary family consisting of the
empty set and all nonempty state sets having a common legal command.
A hyperedge is monochromatic under the
joint probe output exactly when all selected probes are constant on it.
Therefore adequacy is precisely the asserted hitting condition, with the
same declared additive cost.\hfill$\square$

\begin{proposition}[No coarsest online-safe repair]
\label{online-no-coarsest}
Unlike exact autonomous behavior preservation, command-safe observation can
have incomparable minimal adequate refinements and no coarsest one.
\end{proposition}

\emph{Proof.}
Take $E_0=\{a,b\}$, $E_1=\{b,c\}$, $E_2=\{a,c\}$. Every pair shares a
command, but the triple does not. Thus $\mathcal H=\{\{0,1,2\}\}$: pairwise
conflict graphs miss the obstruction. Both partitions
$\{0,1\}\mid\{2\}$ and $\{0\}\mid\{1,2\}$ are adequate. Any common
coarsening identifies $0\sim1\sim2$, so is constant and inadequate. No
coarsest adequate refinement exists.\hfill$\square$

This hypergraph mechanism is classical. It prevents equating the coarsest
autonomous quotient with a least online-capability repair. The theorem is
about supplied fixed probes, not adaptive destructive acquisition.

\paragraph{Verification and authority.}
The accompanying exact Python solver emits positive strategy and negative
rank tables checked by a separate verifier. All 256 two-state/two-command
arenas and 60 further three-state/two-request arenas agree with a distinct
positional-strategy enumeration oracle. A separate public-policy tree oracle
over complete private histories checks 256 public command tables and 768
initial beliefs. These tests validate implementations; the proofs establish
the quantified statements. Generic certificate soundness has a separately
kernel-checked Lean development in \texttt{ForgeGameCertificates.lean}; the
fixed-point algorithm and public-belief compiler are not formalized there.
No academic priority or physical source authentication is asserted.


\section{Authority, prior art, and remaining questions}
The stable quotient and coding arguments are classical. Adaptive
identification and its failures belong to established finite-state testing
\cite{Gunicen,Hierons}; in particular, incomplete distinguishing-sequence
covering is close prior art to the feasible-block formulation. Classical
original/current state bookkeeping is explicit in that work's Lemma 6.2 and
Section 7.1; it is not a new source-indexing invention. Minimum-cost adaptive
identification is also established research \cite{TurkerCost}. Classical
zero-error coding connects information with coloring \cite{Witsenhausen}.
The present forbidden hyperedges may have size three with no forbidden
pairs, which is why a pair-confusability graph alone is insufficient here.
Academic novelty of the exact combined cost/source-memory formulation remains
unverified. This manuscript should not be described as a new general
automata or information theory.

The implementation uses exact rational arithmetic and checks local cost
inequalities plus termination ranks. Tests include 160 fresh four-state
repair models compared with exhaustive partitions and 90 fresh two-original
probe models compared with independent enumeration of stationary policies
simulating each fixed original. A complete 13-belief certificate accompanies
the three-source example. A separate raw-table oracle rebuilds beliefs and
checks all 6,561 two-carrier probe tables plus 140 fresh three-origin models,
including 482 advice-budget comparisons. These are finite implementation checks, not
replacement proofs or external source authentication.

The highest-value extensions are a total finite-controller memory bound and
a source-preserving restoration requirement after acquisition. Both modify
the feasible family and require new arguments. Unbounded carriers,
infinite completion, richer proof authority, and physical access remain
outside the present theorems.

\begin{thebibliography}{9}
\bibitem{Moore} E. F. Moore, Gedanken-experiments on sequential machines,
\emph{Automata Studies}, Princeton University Press, 1956, pp. 129--153.
\bibitem{Rutten} J. J. M. M. Rutten, Universal coalgebra: a theory of systems,
\emph{Theoretical Computer Science} 249 (2000), 3--80.
\bibitem{Gunicen} C. G\"uni\c cen, G.-V. Jourdan and H. Yenig\"un,
Using multiple adaptive distinguishing sequences for checking sequence
generation, author manuscript, \url{https://www.site.uottawa.ca/~gvj/papers/C52.pdf}.
\bibitem{Hierons} R. M. Hierons and U. C. T\"urker, Incomplete distinguishing
sequences for finite state machines, \emph{The Computer Journal} 58 (2015),
3089--3113, doi:10.1093/comjnl/bxv041.
\bibitem{Witsenhausen} H. S. Witsenhausen, The zero-error side information
problem and chromatic numbers, \emph{IEEE Transactions on Information Theory}
22 (1976), 592--593.
\bibitem{TurkerCost} U. C. T\"urker, T. \"Unl\"uyurt and H. Yenig\"un,
Effective algorithms for constructing minimum cost adaptive distinguishing
sequences, \emph{Information and Software Technology} 74 (2016), 69--85,
doi:10.1016/j.infsof.2016.02.001.
\end{thebibliography}
\end{document}
